% TOP_PROBLEM: {"id": 2580, "title": "Kourovka Notebook Problem 21.71", "category_id": 17, "category": {"id": 17, "name": "group_theory", "display_name": "Group Theory", "description": "Problems about groups, group actions, representations, and related algebraic structures.", "slug": "group-theory", "order_index": 17, "created_at": "2026-07-31T15:26:25.671Z"}, "status": "open", "problem_number": "KOU-21.71", "set_id": 10}
% FIRST_POSTED: 2026-10-01
% ENTRY_KIND: statement_only
% UnsolvedMath ID 2580; source catalogue code KOU-21.71
% !TeX program = lualatex
% Definition and sources only; this file is not an LLM solution attempt.
\documentclass[11pt,a4paper]{article}
\usepackage[margin=25mm]{geometry}
\usepackage{fontspec}
\setmainfont{lmroman10-regular.otf}[BoldFont=lmroman10-bold.otf,
 ItalicFont=lmroman10-italic.otf,BoldItalicFont=lmroman10-bolditalic.otf]
\usepackage{amsmath,amssymb,mathtools}
\usepackage{unicode-math}
\setmathfont{latinmodern-math.otf}
\AtBeginDocument{\let\setminus\smallsetminus}
\usepackage{xurl,hyperref}
\hypersetup{colorlinks=true,urlcolor=blue}
\setcounter{secnumdepth}{0}
\setlength{\parindent}{0pt}
\setlength{\parskip}{5pt}
\setlength{\emergencystretch}{3em}
\begin{document}
\section*{Kourovka Notebook Problem 21.71}\label{2580}
{\small UnsolvedMath ID 2580 · KOU-21.71 · Group Theory · Kourovka Notebook - New Problems, Issue 21}
\subsection{Definitions and mathematical statement}
Let $F$ be a field and $FG$ the group algebra of an abstract group
$G$. Give $F$ the trivial $FG$-module structure, so every $g\in G$
acts as the identity. The group has type $FL(F)$ if there is an exact
sequence
\[
 0\longrightarrow P_d\longrightarrow\cdots\longrightarrow P_0
 \longrightarrow F\longrightarrow0
\]
for some finite $d$, with each $P_i$ a finitely generated free
$FG$-module. The differentials are $FG$-linear and free module ranks
are finite; finite generation of the underlying $F$-vector spaces
$P_i$ is not required.

Tensor such a resolution over $FG$ with the trivial right module $F$.
The resulting finite complex of finite-dimensional vector spaces has
homology $H_i(G;F)$, the group homology with trivial coefficients.
Consequently its Euler characteristic is the well-defined integer
\[
 \chi_F(G)=\sum_{i\ge0}(-1)^i\dim_F H_i(G;F),
\]
with only finitely many nonzero summands. It also equals the alternating
sum of the free module ranks in the displayed resolution.

Do there exist a group $G$ and two fields $F_1,F_2$ such that $G$ is
of type $FL(F_1)$ and type $FL(F_2)$, but
$\chi_{F_1}(G)\ne\chi_{F_2}(G)$?
The same abstract group must serve for both coefficient fields.
The two free resolutions may differ in length, ranks and maps.
No finite classifying space, integral free resolution, or finite
presentation is included in the hypothesis; any of those would be
an additional condition beyond the stated problem.
\subsection{Short English statement}
Can a group admitting finite free resolutions over two fields have different Euler characteristics over those fields?
\subsection{Sources}
\begin{itemize}
\item[S1] ULAM.AI and the UnsolvedMath Contributors, supplied problems.json, record 2580 (KOU-21.71), Kourovka Notebook - New Problems, Issue 21. Catalogue identity and source transcription; CC BY 4.0.
\url{https://huggingface.co/datasets/ulamai/UnsolvedMath}.
\item[S2] The Kourovka Notebook, 21st edition, new problems, Problem 21.71. Expanded formulation of the supplied catalogue statement.
\url{https://arxiv.org/pdf/1401.0300v47}.
\end{itemize}
\end{document}
